% ============================================================
\section{End (Endomaps)}
A categoria de endomapas (ou sistemas dinâmicos discretos autónomos) foca-se na evolução de estados.
Um objeto é um conjunto finito munido de uma função (endomorfismo) que mapeia o conjunto para si mesmo, $\alpha: X \to X$.

\subsection{Objetos}
Dado que a relação é estritamente uma função (cada elemento tem exatamente uma imagem), a representação mais eficiente em MATLAB não é uma matriz, mas sim um vetor de índices, o que otimiza significativamente o uso de memória e tempo de acesso.

\begin{lstlisting}[caption={Representação de um Endomapa}]
% Vetor representando f(1)=2, f(2)=3, f(3)=1 (Um ciclo de comprimento 3)
E = [2, 3, 1];
\end{lstlisting}

\subsection{Morfismos}
Os morfismos são funções equivariantes que "comutam" com a dinâmica.
Um mapeamento $f: (X, \alpha) \to (Y, \beta)$ é válido se aplicar a dinâmica e depois mapear for igual a mapear e depois aplicar a dinâmica do codomínio: $f(\alpha(x)) = \beta(f(x))$.

\begin{lstlisting}[caption={Verificação do Quadrado Comutativo}]
% f é o morfismo; EA e EB são os endomapas
isMorphism = true;
for x = 1:length(f)
    % Verifica se comuta: f(alpha(x)) == beta(f(x))
    if f(EA(x)) ~= EB(f(x))
        isMorphism = false;
        break;
    end
end
\end{lstlisting}